\documentclass[12pt]{jarticle}
\usepackage[utf8]{inputenc}
\usepackage[top=30truemm, bottom=30truemm, left=20truemm, right=20truemm]{geometry}
\usepackage{amsmath}
\usepackage{amssymb}
\usepackage{mathtools}
\usepackage{siunitx}
\usepackage{bm}
\usepackage[dvipdfmx]{graphicx}
\usepackage[dvipdfmx]{color}
\usepackage[subrefformat=parens]{subcaption}
\usepackage{booktabs}
\usepackage{tabularx}
\usepackage{tocloft}
\usepackage{enumerate}
\usepackage{url}
\usepackage{multirow}
\usepackage{caption}
\usepackage{enumitem}
\usepackage{makecell}
\usepackage{array}
\usepackage{algorithm}
\usepackage{algpseudocode}
\usepackage{amsthm}
\usepackage{physics}
\usepackage{mathrsfs}
\usepackage{bbm}

\numberwithin{equation}{section}    % 数式番号にセクション番号をつける
\numberwithin{figure}{section}      % 図番号にセクション
\numberwithin{table}{section}      % 図番号にセクション

% \renewcommand{\figurename}{Fig.}    % 図 -> Fig.
% \renewcommand{\tablename}{Table }   % 表 -> Table

\renewcommand{\baselinestretch}{1.1}

\theoremstyle{definition}
\newtheorem{theorem}{Theorem}
\newtheorem{proposition}{Proposition}
\newtheorem{lemma}{Lemma}
\newtheorem{definition}{Definition}

\newlength{\figcaptionskip}
\setlength{\figcaptionskip}{5pt} % 図のキャプション間隔
\newlength{\tabcaptionskip}
\setlength{\tabcaptionskip}{-5pt} % 表のキャプション間隔
\captionsetup[figure]{skip=\figcaptionskip}
\captionsetup[table]{skip=\tabcaptionskip}

\setlist[enumerate]{topsep=5pt, partopsep=5pt, itemsep=0pt, parsep=0pt}
\setlength{\topsep}{5pt}
\setlength{\partopsep}{5pt}

\DeclareMathOperator*{\argmin}{arg\,min}

\begin{document}

\tableofcontents
\newpage

\section{比率の差の検定}

\subsection{どのような検定か}

比率の差の検定は、ABテストにおいて、A群とB群の間でCVRに有意な差があるかを判定する場合に用いることができる。二項分布に従う確率変数$X_{1} \sim B(n_{1}, p_{1}),\ X_{2} \sim B(n_{2}, p_{2})$における母比率$p_{1}$と$p_{2}$に違いがあるかを判断するための検定である。ここで、$n_{1}, n_{2}$は各群のサンプルサイズであり、$n_{2} = rn_{1}$とする。また、
\begin{equation}
    \quad \hat{p}_{1} = \frac{X_{1}}{n_{1}}, \quad \hat{p}_{2} = \frac{X_{2}}{n_{2}}
\end{equation}
とする。いま、$n_{1}, n_{2}$は十分大きく、中心極限定理より
\begin{equation}
    \hat{p}_{1} \sim N\left(p_{1}, \frac{p_{1}(1 - p_{1})}{n_{1}}\right), \quad \hat{p}_{2} \sim N\left(p_{2}, \frac{p_{2}(1 - p_{2})}{n_{2}}\right)
\end{equation}
が成り立つとする。これは、$X_{k}\ (k = 1, 2)$が$B(1, p_{k})$に従う確率変数の和であることによる。このとき、$\hat{\delta} = \hat{p}_{1} - \hat{p}_{2}$は正規分布の再生性より
\begin{equation}
    \hat{\delta} \sim N\left(p_{1} - p_{2},\ \frac{p_{1}(1 - p_{1})}{n_{1}} + \frac{p_{2}(1 - p_{2})}{n_{2}}\right)
\end{equation}
となる。特に、帰無仮説$H_{0}: p_{1} = p_{2} (= p)$が成り立つ時、
\begin{equation}
    \hat{\delta} \sim N\left(0,\ p(1 - p)\left(\frac{1}{n_{1}} + \frac{1}{n_{2}}\right)\right)
\end{equation}
である。$\hat{p} = (X_{1} + X_{2}) / (n_{1} + n_{2})$として、
\begin{equation}
    Z = \frac{\hat{\delta}}{\sqrt{\hat{p}(1 - \hat{p})\left(\frac{1}{n_{1}} + \frac{1}{n_{2}}\right)}}
\end{equation}
が検定統計量である。対立仮説に対する棄却域とp値の計算方法は表\ref{sec1:tab:rejection_region}の通りである。$z$は正規分布に従う確率変数であり、$z_{\alpha}$は上側100$\alpha$\%を満たす点である。

\begin{table*}[h]
    \centering
    \caption{比率の差の検定における対立仮説と棄却域、p値の計算方法}
    \label{sec1:tab:rejection_region}
    \begin{center}
        \renewcommand{\arraystretch}{0.9} % 行の高さ調整
        \setlength{\tabcolsep}{8pt}      % 列の幅調整
        \scalebox{1.0}{
            \begin{tabular}{|c|c|c|}
                \hline
                対立仮説               & 棄却域                  & p値の計算方法       \\
                \hline
                $p_{1} \neq p_{2}$ & $|Z| > z_{\alpha/2}$ & $2P(z > |Z|)$ \\
                $p_{1} > p_{2}$    & $Z > z_{\alpha}$     & $P(z > Z)$    \\
                $p_{1} < p_{2}$    & $Z < -z_{\alpha}$    & $P(z < Z)$    \\
                \hline
            \end{tabular}
        }
    \end{center}
\end{table*}

\subsection{サンプルサイズの設計方法}

\subsubsection{両側検定の場合}
対立仮説を$p_{1} \neq p_{2}$としたときのサンプルサイズ設計に用いる計算式を導出する。対立仮説が成り立つとした時、
\begin{equation}
    1 - \beta = P(|Z| > z_{\alpha/2}) = P(Z > z_{\alpha/2}) + P(Z < -z_{\alpha/2}) \label{sec1:eq:stat_power}
\end{equation}
である。ここで、$p_{1} > p_{2}$とする。この時、式\eqref{sec1:eq:stat_power}の第二項は十分小さいとして無視する。すると、以下のように整理できる。
\begin{align}
    1 - \beta & = P(Z > z_{\alpha/2}) \label{sec1:eq:stat_power_simplified}                                                                                                                                                                                                                                                            \\
              & = P\left(\frac{\hat{\delta}}{\sqrt{p(1 - p)\left(\frac{1}{n_{1}} + \frac{1}{n_{2}}\right)}} > z_{\alpha/2}\right)                                                                                                                                                                                                      \\
              & = P\left(\frac{\hat{\delta} - (p_{1} - p_{2})}{\sqrt{\frac{p_{1}(1 - p_{1})}{n_{1}} + \frac{p_{2}(1 - p_{2})}{n_{2}}}} > \frac{z_{\alpha/2}\sqrt{\hat{p}(1 - \hat{p})\left(\frac{1}{n_{1}} + \frac{1}{n_{2}}\right)} - (p_{1} - p_{2})}{\sqrt{\frac{p_{1}(1 - p_{1})}{n_{1}} + \frac{p_{2}(1 - p_{2})}{n_{2}}}}\right) \\
\end{align}
対立仮説が真である時に
\begin{equation}
    \frac{\hat{\delta} - (p_{1} - p_{2})}{\sqrt{\frac{p_{1}(1 - p_{1})}{n_{1}} + \frac{p_{2}(1 - p_{2})}{n_{2}}}} \sim N(0, 1)
\end{equation}
であるから、
\begin{equation}
    -z_{\beta} = \frac{z_{\alpha/2}\sqrt{\hat{p}(1 - \hat{p})\left(\frac{1}{n_{1}} + \frac{1}{n_{2}}\right)} - (p_{1} - p_{2})}{\sqrt{\frac{p_{1}(1 - p_{1})}{n_{1}} + \frac{p_{2}(1 - p_{2})}{n_{2}}}}
\end{equation}
が成り立つ。ここで、$n_{2} = rn_{1}$を使って$n_{1}$について解くと、
\begin{equation}
    n_{1} = \frac{\left(z_{\alpha/2}\sqrt{\hat{p}(1 - \hat{p})(r + 1)} + z_{\beta}\sqrt{rp_{1}(1 - p_{1}) + p_{2}(1 - p_{2})}\right)^{2}}{(p_{1} - p_{2})^{2}r}
\end{equation}
を得る。$\hat{p} = (X_{1} + X_{2})/(n_{1} + n_{2})$であったが、以下のように整理すれば$n_{1}, n_{2}$に依存せず計算可能である。
\begin{align}
    \hat{p} & = \frac{X_{1} + X_{2}}{n_{1} + n_{2}} = \frac{n_{1}\hat{p_{1}} + n_{2}\hat{p}_{2}}{n_{1} + n_{2}} = \frac{n_{1}\hat{p_{1}} + rn_{1}\hat{p}_{2}}{n_{1} + rn_{1}} = \frac{\hat{p}_{1} + r\hat{p}_{2}}{1 + r} \\
\end{align}

上記の導出は$p_{1} > p_{2}$を仮定したが、$p_{1} < p_{2}$を仮定した場合は
\begin{equation}
    1 - \beta = P(Z < -z_{\alpha/2}) = P(Z > z_{\alpha/2})
\end{equation}
となるため、同じ計算式で良い。

\subsubsection{片側検定の場合}
対立仮説を$p_{1} > p_{2}$とした場合のサンプルサイズ設計に用いる計算式は、両側検定における計算式の$z_{\alpha/2}$に$z_{\alpha}$を代入すれば良い。式\eqref{sec1:eq:stat_power_simplified}が変わるだけだからである。

% \subsection{カイ二乗検定と比率の差の検定の関係}
% 次に、表\ref{tab:crosstab}のような$2 \times 2$のクロス集計表に対して用いられるカイ二乗検定を考える。
% \begin{table}[hbt]
%     \centering
%     \caption{観測データのクロス集計表}
%     \label{tab:crosstab}
%     \vspace{5pt}
%     \begin{tabular}{|c|cc|c|}
%         \hline
%                         & $B_{1}$       & $B_{2}$       & 合計$T_{i \cdot}$ \\
%         \hline
%         $A_{1}$         & $x_{11}$      & $x_{12}$      & $T_{1 \cdot}$   \\
%         $A_{2}$         & $x_{21}$      & $x_{22}$      & $T_{2 \cdot}$   \\
%         \hline
%         合計$T_{\cdot j}$ & $T_{\cdot 1}$ & $T_{\cdot 2}$ & $T$             \\
%         \hline
%     \end{tabular}
% \end{table}
% 確率を表\ref{tab:crosstab_prob}のようにおく。
% \begin{table}[hbt]
%     \centering
%     \caption{一様性の検定のための確率のおき方}
%     \label{tab:crosstab_prob}
%     \vspace{5pt}
%     \begin{tabular}{|c|cc|c|}
%         \hline
%                 & $B_{1}$  & $B_{2}$  & 合計 \\
%         \hline
%         $A_{1}$ & $p_{11}$ & $p_{12}$ & 1  \\
%         $A_{2}$ & $p_{21}$ & $p_{22}$ & 1  \\
%         \hline
%     \end{tabular}
% \end{table}
% この確率を用いて、$A_{1}$群で効果がある確率を$A_{2}$群で効果がある確率が等しいという帰無仮説は$p_{11} = p_{21}$と表される。対立仮説は$p_{11} \neq p_{21}$とする。

% 帰無仮説のもとで、$(A_{i}, B_{j})$の期待度数$t_{ij}$は
% \begin{equation}
%     t_{ij} = \frac{T_{\cdot j}}{T} T_{i \cdot}
% \end{equation}
% で求められる。全体における効果が出る割合$T_{\cdot j} / T$にi群の人数$T_{i \cdot}$をかけている。表\ref{tab:crosstab_expected_freq}に期待度数を示す。
% \begin{table}[hbt]
%     \centering
%     \caption{期待度数}
%     \label{tab:crosstab_expected_freq}
%     \vspace{5pt}
%     \begin{tabular}{|c|cc|c|}
%         \hline
%                         & $B_{1}$       & $B_{2}$       & 合計$T_{i \cdot}$ \\
%         \hline
%         $A_{1}$         & $t_{11}$      & $t_{12}$      & $T_{1 \cdot}$   \\
%         $A_{2}$         & $t_{21}$      & $t_{22}$      & $T_{2 \cdot}$   \\
%         \hline
%         合計$T_{\cdot j}$ & $T_{\cdot 1}$ & $T_{\cdot 2}$ & $T$             \\
%         \hline
%     \end{tabular}
% \end{table}

% この時、
% \begin{equation}
%     \chi^{2} = \sum_{i = 1}^{2}\sum_{j = 1}^{2} \frac{(x_{ij} - t_{ij})^{2}}{t_{ij}}
% \end{equation}
% の確率分布は、帰無仮説のもとで自由度1のカイ二乗分布に近似できる。これより、$\chi^{2}$を検定統計量としたカイ二乗検定が可能である。カイ二乗分布の定義より、$Z^{2}$も自由度1のカイ二乗分布に従うことに注意する。ここで、$\chi^{2}_{\alpha}$をカイ二乗分布における上側$100\alpha$\%点とすると、
% \begin{equation}
%     P(\chi^{2} > \chi^{2}_{\alpha}(1)) = \alpha
% \end{equation}
% かつ
% \begin{equation}
%     P(|Z| > z_{\alpha / 2}) = P(Z^{2} > z_{\alpha / 2}^{2}) = P(\chi^{2} > z_{\alpha / 2}^{2}) = \alpha
% \end{equation}
% であるから、
% \begin{equation}
%     P(\chi^{2} > \chi^{2}_{\alpha}(1)) = P(\chi^{2} > z_{\alpha / 2}^{2})
% \end{equation}
% よって、$\chi_{\alpha}(1) = z_{\alpha / 2}^{2}$が成り立つ。

% ここで、比率の差の検定における検定統計量$Z$は
% \begin{equation}
%     Z = \frac{\hat{\delta}}{\sqrt{(1 + r)p(1 - p) / rn_{1}}} \sim N(0, 1)
% \end{equation}
% であった（実際は$p$に$\overline{p}$を代入する）。また、$2 \times 2$のクロス集計表に対するカイ二乗検定は、比率の差の検定として見ることもできる。なぜなら、互いに独立な二項分布に従う確率変数$X_{11}, X_{21}$
% \begin{equation}
%     X_{11} \sim B(T_{1 \cdot}, p_{11}), \quad X_{21} \sim B(T_{2 \cdot}, p_{21})
% \end{equation}
% における$p_{11}$の$p_{21}$の同等性を扱う検定問題だからである。

% 今、$p_{11}$と$p_{21}$の間で検出したい差を$\epsilon$とし、この検出したい差に対して有意水準$\alpha$で検出力$1 - \beta$を持つ比率の差の検定が行えるサンプルサイズを、\eqref{eq:proportion_sample_size}で計算したとする。この時、帰無仮説のもとでは
% \begin{equation}
%     P(\chi^{2} > \chi^{2}_{\alpha}(1)) = \alpha
% \end{equation}
% であり、対立仮説のもとでは
% \begin{equation}
%     1 - \beta = P(|Z| > z_{\alpha / 2}) = P(Z^{2} > z_{\alpha / 2}^{2}) = P(\chi^{2} > z_{\alpha / 2}^{2}) = P(\chi^{2} > \chi_{\alpha}(1))
% \end{equation}
% である。ここでは、比率の差の検定の検出力が$1 - \beta$となるサンプルサイズであること、$Z^{2} = \chi^{2}$であること、$\chi_{\alpha}(1) = z_{\alpha / 2}^{2}$であることを用いた。

% これより、比率の差の検定のために算出されたサンプルサイズを確保すれば、カイ二乗検定も有意水準$\alpha$で検出力$1 - \beta$を持った検定となることがわかる。


\section{Welchのt検定}

\subsection{どのような検定か}

Welchのt検定は、ABテストにおいて、A群とB群の間でユーザー当たりの平均CV数に有意な差があるかを判定する場合に用いることができる。確率変数$X_{1,i}\ (i=1,\ldots,n_{1}) \sim N(\mu_{1}, \sigma_{1}^{2})$の母平均$\mu_{1}$と、$X_{2, i}\ (i=1,\ldots,n_{2}) \sim N(\mu_{2}, \sigma_{2}^{2})$の母平均$\mu_{2}$が異なるかを判断するための検定である。ここで、$n_{1},n_{2}$は各群のサンプルサイズであり、$n_{2}=rn_{1}$とする。
また、$\mathbb{E}[X_{1,i}]=\mu_{1}, \text{Var}[X_{1,i}]=\sigma_{1}^{2}, \mathbb{E}[X_{2,i}]=\mu_{2}, \text{Var}[X_{2,i}]=\sigma_{2}^{2}$である。

Welchのt検定における検定統計量は次の通りである。
\begin{equation}
    T = \frac{\bar{X}_1 - \bar{X}_2}{\sqrt{\frac{s_{1}^{2}}{n_{1}} + \frac{s_{2}^{2}}{n_{2}}}}
\end{equation}
ここで、
\begin{align}
    \bar{X_{1}} & = \frac{1}{n} \sum_{i = 1}^{n} X_{1, i}                         \\
    \bar{X_{2}} & = \frac{1}{n} \sum_{i = 1}^{n} X_{2, i}                         \\
    s_{1}^{2}   & = \frac{1}{n - 1} \sum_{i = 1}^{n} (X_{1, i} - \bar{X}_{1})^{2} \\
    s_{2}^{2}   & = \frac{1}{n - 1} \sum_{i = 1}^{n} (X_{2, i} - \bar{X}_{2})^{2}
\end{align}
である。対立仮説に対する棄却域とp値の計算方法は表\ref{sec2:tab:rejection_region}の通りである。$t$はt分布に従う確率変数であり、$t_{\alpha}(\nu)$は自由度$\nu$のt分布において上側100$\alpha$\%を満たす点である。

\begin{table*}[h]
    \centering
    \caption{t検定における対立仮説と棄却域、p値の計算方法}
    \label{sec2:tab:rejection_region}
    \begin{center}
        \renewcommand{\arraystretch}{0.9} % 行の高さ調整
        \setlength{\tabcolsep}{8pt}      % 列の幅調整
        \scalebox{1.0}{
            \begin{tabular}{|c|c|c|}
                \hline
                対立仮説                   & 棄却域                       & p値の計算方法       \\
                \hline
                $\mu_{1} \neq \mu_{2}$ & $|T| > t_{\alpha/2}(\nu)$ & $2P(t > |T|)$ \\
                $\mu_{1} > \mu_{2}$    & $T > t_{\alpha}(\nu)$     & $P(t > T)$    \\
                $\mu_{1} < \mu_{2}$    & $T < -t_{\alpha}(\nu)$    & $P(t < T)$    \\
                \hline
            \end{tabular}
        }
    \end{center}
\end{table*}

\subsection{サンプルサイズの設計}

\subsubsection{両側検定の場合}

対立仮説を$\mu_{1} \neq \mu_{2}$としたときのサンプルサイズ設計に用いる計算式を導出する。対立仮説が成り立つとした時、
\begin{equation}
    1 - \beta = P(|T| > t_{\alpha/2}) = P(T > t_{\alpha/2}) + P(T < -t_{\alpha/2}) \label{sec2:eq:stat_power}
\end{equation}
である。ここで、$\mu_{1} > \mu_{2}$とする。この時、式\eqref{sec2:eq:stat_power}の第二項は十分小さいとして無視する。また、サンプルサイズが大きい時に検定統計量は標準正規分布に従う確率変数に分布収束することを利用し、以下のように整理する。
\begin{align}
    1 - \beta & = P(T > z_{\alpha/2})                                                                                                                                                                                                                \\
              & = P\left(\frac{\bar{X}_{1} - \bar{X}_{2}}{\sqrt{\frac{s_{1}^{2}}{n_{1}} + \frac{s_{2}^{2}}{n_{2}}}} > z_{\alpha/2}\right)                                                                                                            \\
              & = P\left(\frac{\bar{X}_{1} - \bar{X}_{2} - (\mu_{1} - \mu_{2})}{\sqrt{\frac{s_{1}^{2}}{n_{1}} + \frac{s_{2}^{2}}{n_{2}}}} > z_{\alpha/2} - \frac{\mu_{1} - \mu_{2}}{\sqrt{\frac{s_{1}^{2}}{n_{1}} + \frac{s_{2}^{2}}{n_{2}}}}\right)
\end{align}
これより、
\begin{align}
    -z_{\beta} = z_{\alpha/2} - \frac{\mu_{1} - \mu_{2}}{\sqrt{\frac{s_{1}^{2}}{n_{1}} + \frac{s_{2}^{2}}{n_{2}}}}
\end{align}
が成り立つ。$n_{2} = rn_{1}$を利用して$n_{1}$について解くことで、サンプルサイズの計算式を得る。
\begin{equation}
    n_{1} = \frac{(z_{\alpha/2} + z_{\beta})^{2} \left(s_{1}^{2} + \frac{s_{2}^{2}}{r}\right)}{(\mu_{1} - \mu_{2})^{2}}
\end{equation}

\subsubsection{片側検定の場合}
両側検定の場合の計算式において、$z_{\alpha/2}$に$z_{\alpha}$を代入すれば良い。

\subsection{正規分布に従っていないデータに対する適用}

ABテストにおいて、サンプルサイズが多い場合はデータの性質によらず、Welchのt検定を用いて良いかを検討する。

$X_{1, 1}, \ldots, X_{1, n}$が群1、$X_{2, 1}, \ldots, X_{2, n}$が群2のサンプルであり、それぞれ群内で独立に同一分布に従うとする。また、群1と群2は独立であるとする。$\mathbb{E}[X_{1, i}] = \mu_1, \text{Var}[X_{1, i}] = \sigma^2_1, \mathbb{E}[X_{2, i}] = \mu_2, \text{Var}[X_{2, i}] = \sigma^2_2$とする($\sigma_{1}^{2}, \sigma_{2}^{2} \in (0, \infty)$)。

Welchのt検定における検定統計量は
\begin{equation}
    t = \frac{\bar{X}_1 - \bar{X}_2 - (\mu_1 - \mu_2)}{\sqrt{\frac{s_{1}^{2} + s_{2}^{2}}{n}}}
\end{equation}
である（帰無仮説が真であることは仮定しない）。ここで、
\begin{align}
    \bar{X_{1}} & = \frac{1}{n} \sum_{i = 1}^{n} X_{1, i}                         \\
    \bar{X_{2}} & = \frac{1}{n} \sum_{i = 1}^{n} X_{2, i}                         \\
    s_{1}^{2}   & = \frac{1}{n - 1} \sum_{i = 1}^{n} (X_{1, i} - \bar{X}_{1})^{2} \\
    s_{2}^{2}   & = \frac{1}{n - 1} \sum_{i = 1}^{n} (X_{2, i} - \bar{X}_{2})^{2}
\end{align}
である。また、
\begin{align}
    U_{n} & = \frac{\bar{X}_1 - \bar{X}_2 - (\mu_1 - \mu_2)}{\sqrt{\frac{\sigma_{1}^{2} + \sigma_{2}^{2}}{n}}} \\
    R_{n} & = \sqrt{\frac{s_{1}^{2} + s_{2}^{2}}{\sigma_{1}^{2} + \sigma_{2}^{2}}}
\end{align}
とする。

$U_{n}$の分子は、
\begin{equation}
    \bar{X}_1 - \bar{X}_2 - (\mu_1 - \mu_2) = \frac{1}{n} \sum_{i = 1}^{n} ((X_{1, i} - \mu_{1}) - (X_{2, i} - \mu_{2}))
\end{equation}
と整理できる。ここで、$W_{i} = (X_{1, i} - \mu_{1}) - (X_{2, i} - \mu_{2})$とおく。群1内のi.i.d.性、群2内のi.i.d.性、および群間独立性より、$W_{1}, \ldots, W_{n}$もi.i.d.である。$W_{i}$の期待値と分散はそれぞれ次のようになる。
\begin{align}
    \mathbb{E}[W_{i}] & = \mathbb{E}[(X_{1, i} - \mu_{1}) - (X_{2, i} - \mu_{2})]         \\
                      & = \mathbb{E}[X_{1, i}] - \mu_{1} - \mathbb{E}[X_{2, i}] + \mu_{2} \\
                      & = 0
\end{align}
$X_{1, i}$と$X_{2, i}$の独立性より、
\begin{align}
    \text{Var}[W_{i}] & = \text{Var}[(X_{1, i} - \mu_{1}) - (X_{2, i} - \mu_{2})] \\
                      & = \text{Var}[X_{1, i}] + \text{Var}[X_{2, i}]             \\
                      & = \sigma_{1}^{2} + \sigma_{2}^{2}
\end{align}
これは正の有限値である。中心極限定理より、
\begin{equation}
    U_{n} = \frac{\bar{W}}{\sqrt{\frac{\sigma_{1}^{2} + \sigma_{2}^{2}}{n}}} = \frac{\sqrt{n}(\bar{W} - 0)}{\sqrt{\sigma_{1}^{2} + \sigma_{2}^{2}}} \xrightarrow{d} Z \sim \mathcal{N}(0, 1)
\end{equation}
となる。

一方、$R_{n}$について、$s_{1}^{2} \xrightarrow{p} \sigma_{1}^{2}, s_{2}^{2} \xrightarrow{p} \sigma_{2}^{2}$が成り立つことを示す。$k = 1, 2$について、
\begin{align}
    \sum_{i = 1}^{n} (X_{k, i} - \bar{X}_{k})^{2}
     & = \sum_{i = 1}^{n} ((X_{k, i} - \mu_{k}) - (\bar{X}_{k} - \mu_{k}))^{2}                                                                     \\
     & = \sum_{i = 1}^{n} (X_{k, i} - \mu_{k})^{2} - 2(\bar{X}_{k} - \mu_{k}) \sum_{i = 1}^{n} (X_{k, i} - \mu_{k}) + n(\bar{X}_{k} - \mu_{k})^{2}
\end{align}
となる。ここで、$\sum_{i = 1}^{n} (X_{k, i} - \mu_{k}) = n(\bar{X}_{k} - \mu_{k})$より、
\begin{align}
    \sum_{i = 1}^{n} (X_{k, i} - \bar{X}_{k})^{2}
     & = \sum_{i = 1}^{n} (X_{k, i} - \mu_{k})^{2} - 2n(\bar{X}_{k} - \mu_{k})^{2} + n(\bar{X}_{k} - \mu_{k})^{2} \\
     & = \sum_{i = 1}^{n} (X_{k, i} - \mu_{k})^{2} - n(\bar{X}_{k} - \mu_{k})^{2}
\end{align}
となる。よって、
\begin{equation}
    s_{k}^{2} = \frac{n}{n - 1} \left(\frac{1}{n} \sum_{i = 1}^{n} (X_{k, i} - \mu_{k})^{2} - (\bar{X}_{k} - \mu_{k})^{2}\right)
\end{equation}
と整理できる。

大数の弱法則より$(\bar{X}_{k} - \mu_{k}) \xrightarrow{p} 0$であり、連続写像定理から$(\bar{X}_{k} - \mu_{k})^{2} \xrightarrow{p} 0$である。また、
\begin{equation}
    \mathbb{E}\left[\frac{1}{n} \sum_{i = 1}^{n} (X_{k, i} - \mu_{k})^{2}\right] = \mathbb{E}[(X_{k} - \mu_{k})^{2}] = \text{Var}[X_{k}] = \sigma_{k}^{2}
\end{equation}
となるから、大数の弱法則より$\frac{1}{n} \sum_{i = 1}^{n} (X_{k, i} - \mu_{k})^{2} \xrightarrow{p} \sigma_{k}^{2}$である。

また、$n \to \infty$のとき$\frac{n}{n - 1} \to 1$であるから、連続写像定理より、
\begin{equation}
    s_{k}^{2} = \frac{n}{n - 1} \left(\frac{1}{n} \sum_{i = 1}^{n} (X_{k, i} - \mu_{k})^{2} - (\bar{X}_{k} - \mu_{k})^{2}\right) \xrightarrow{p} \sigma_{k}^{2}
\end{equation}
である。

ここで、「統計学への確率論、その先へ」の定理4.3.1より、$(s_{1}^{2}, s_{2}^{2}) \xrightarrow{p} (\sigma_{1}^{2}, \sigma_{2}^{2})$が成り立つ。再び連続写像定理より
\begin{equation}
    R_{n} = \sqrt{\frac{s_{1}^{2} + s_{2}^{2}}{\sigma_{1}^{2} + \sigma_{2}^{2}}} \xrightarrow{p} \sqrt{\frac{\sigma_{1}^{2} + \sigma_{2}^{2}}{\sigma_{1}^{2} + \sigma_{2}^{2}}} = 1
\end{equation}
となる。

スラツキーの定理を用いて
\begin{equation}
    t = \frac{U_{n}}{R_{n}} \xrightarrow{d} Z
\end{equation}
となる。

これで、元のデータが正規分布に従っていなかったとしても、サンプルサイズが十分大きいのであれば、検定統計量は標準正規分布に従う確率変数に分布収束するとわかった。そのため、棄却域を与えるパーセント点は標準正規分布から与えれば良い。また、サンプルサイズが十分大きいのであれば、Welch-Satterthwaiteの自由度も大きくなり、t分布は標準正規分布で近似できる。これより、データが正規分布に従っていないがサンプルサイズが大きい時、特別に標準正規分布からパーセント点を求めずとも、t分布から求めたパーセント点を用いれば良い。すなわち、Welchのt検定をそのまま用いることができる。


\section{ポアソン分布の差の検定}

\subsection{どのような検定か}
ある事象の発生回数がポアソン分布に従うとし、各群で独立に同一の分布に従う確率変数を$X_{1, i} \sim \text{Po}(\lambda_{1})\ (i = 1, \ldots, n_{1}),\ X_{2, i} \sim \text{Po}(\lambda_{2})\ (i = 1, \ldots, n_{2})$とする。このようなデータとして、例えばユーザーあたりのCV数が考えられる。
\begin{equation}
    \hat{\lambda}_{1} = \frac{1}{n_{1}} \sum_{i = 1}^{n_{1}} X_{1, i},\ \hat{\lambda}_{2} = \frac{1}{n_{2}} \sum_{i = 1}^{n_{2}} X_{2, i}
\end{equation}
とすると、$\mathbb{E}[X_{k, i}] = \lambda_{k},\ \text{Var}[X_{k, i}] = \lambda_{k}$より、
\begin{align}
    \mathbb{E}[\hat{\lambda}_{k}] = \lambda_{k},\ \text{Var}[\hat{\lambda}_{k}] = \frac{\lambda_{k}}{n_{k}}
\end{align}
となる。サンプルサイズが十分大きい場合、中心極限定理より$\hat{\lambda_{k}} \sim N(\lambda_{k}, \frac{\lambda_{k}}{n_{k}})$が成り立ち、正規分布の再生性より$\hat{\lambda_{1}} - \hat{\lambda_{2}} \sim N(\lambda_{1} - \lambda_{2}, \frac{\lambda_{1}}{n_{1}} + \frac{\lambda_{2}}{n_{2}})$となる。

以上のことから、帰無仮説$H_{0}$を$\lambda_{1} = \lambda_{2}$とする仮説検定において、検定統計量は
\begin{equation}
    Z = \frac{\hat{\lambda_{1}} - \hat{\lambda_{2}}}{\sqrt{\frac{\hat{\lambda_{1}}}{n_{1}} + \frac{\hat{\lambda_{2}}}{n_{2}}}}
\end{equation}
となる。対立仮説に対する棄却域とp値の計算方法は表\ref{sec3:tab:rejection_region}の通りである。$z$は正規分布に従う確率変数であり、$z_{\alpha}$は上側100$\alpha$\%を満たす点である。

\begin{table*}[h]
    \centering
    \caption{ポアソン分布の差の検定における対立仮説と棄却域、p値の計算方法}
    \label{sec3:tab:rejection_region}
    \begin{center}
        \renewcommand{\arraystretch}{0.9} % 行の高さ調整
        \setlength{\tabcolsep}{8pt}      % 列の幅調整
        \scalebox{1.0}{
            \begin{tabular}{|c|c|c|}
                \hline
                対立仮説                           & 棄却域                  & p値の計算方法       \\
                \hline
                $\lambda_{1} \neq \lambda_{2}$ & $|Z| > z_{\alpha/2}$ & $2P(z > |Z|)$ \\
                $\lambda_{1} > \lambda_{2}$    & $Z > z_{\alpha}$     & $P(z > Z)$    \\
                $\lambda_{1} < \lambda_{2}$    & $Z < -z_{\alpha}$    & $P(z < Z)$    \\
                \hline
            \end{tabular}
        }
    \end{center}
\end{table*}

\subsection{サンプルサイズの設計}

\subsubsection{両側検定}

対立仮説を$\lambda_{1} \neq \lambda_{2}$としたときのサンプルサイズ計算に用いる計算式を導出する。そのほかの検定と同様にして、
\begin{align}
    1 - \beta & = P(Z > z_{\alpha/2})                                                                                                                                                                                                                                                \\
              & = P\left(\frac{\hat{\lambda_{1}} - \hat{\lambda_{2}}}{\sqrt{\frac{\hat{\lambda_{1}}}{n_{1}} + \frac{\hat{\lambda_{2}}}{n_{2}}}} > z_{\alpha/2}\right)                                                                                                                \\
              & = P\left(\frac{\hat{\lambda_{1}} - \hat{\lambda_{2}} - (\mu_{1} - \mu_{2})}{\sqrt{\frac{\hat{\lambda_{1}}}{n_{1}} + \frac{\hat{\lambda_{2}}}{n_{2}}}} > z_{\alpha/2} - \frac{\mu_{1} - \mu_{2}}{\sqrt{\frac{\lambda_{1}}{n_{1}} + \frac{\lambda_{2}}{n_{2}}}}\right)
\end{align}
これより、
\begin{equation}
    -z_{\beta} = z_{\alpha/2} - \frac{\mu_{1} - \mu_{2}}{\sqrt{\frac{\lambda_{1}}{n_{1}} + \frac{\lambda_{2}}{n_{2}}}}
\end{equation}
が成り立つ。これを$n_{1}$について解くことで、サンプルサイズの計算式を得る。
\begin{equation}
    n_{1} = \frac{(z_{\alpha/2} + z_{\beta})^{2}\left(\lambda_{1} + \frac{\lambda_{2}}{r}\right)}{(\mu_{1} - \mu_{2})^{2}}
\end{equation}

\subsubsection{片側検定}
両側検定の計算式における$z_{\alpha/2}$に$z_{\alpha}$を代入すれば良い。


\section{同等性検定}

\subsection{どのような検定か}

同等性検定は、2つの群に同等とみなせる程度の差しかないことを示すための検定である。2群の母平均$\mu_{1}, \mu_{2}$について、同等とみなせる差を$\triangle_{l}, \triangle_{u}\ (\triangle_{l} < \triangle_{u})$とするとき、帰無仮説$H_{0}$と対立仮説$H_{1}$は
\begin{gather}
    H_{0}: \mu_{1} - \mu_{2} < \triangle_{l} \lor \mu_{1} - \mu_{2} > \triangle_{u} \\
    H_{1}: \triangle_{l} \le \mu_{1} - \mu_{2} \le \triangle_{u}
\end{gather}
とする。帰無仮説を棄却するとき、同等とみなせる程度しか差がなかったと言える。

Welchのt検定を用いることができ、検定統計量は
\begin{equation}
    T_{l} = \frac{\bar{X}_{1} - \bar{X}_{2} - \triangle_{l}}{\sqrt{\frac{s_{1}^{2}}{n_{1}} + \frac{s_{2}^{2}}{n_{2}}}},\ T_{r} = \frac{\bar{X}_{1} - \bar{X}_{2} - \triangle_{u}}{\sqrt{\frac{s_{1}^{2}}{n_{1}} + \frac{s_{2}^{2}}{n_{2}}}}
\end{equation}
である。$\mu_{1} - \mu_{2} < \triangle_{l}$に対する棄却域は$T > t_{\alpha}(\nu)$、$\mu_{1} - \mu_{2} > \triangle_{u}$に対する棄却域は$T < -t_{\alpha}(\nu)$である。

\subsection{サンプルサイズの設計}

Welchのt検定と同様に求めることができる。まず、帰無仮説が$\mu_{1} - \mu_{2} < \triangle_{l}$である場合を導出する。
\begin{align}
    1 - \beta & = P(T > z_{\alpha})                                                                                                                                                                                                                                \\
              & = P\left(\frac{\bar{X}_{1} - \bar{X}_{2} - \triangle_{l}}{\sqrt{\frac{s_{1}^{2}}{n_{1}} + \frac{s_{2}^{2}}{n_{2}}}} > z_{\alpha}\right)                                                                                                            \\
              & = P\left(\frac{\bar{X}_{1} - \bar{X}_{2} - (\mu_{1} - \mu_{2})}{\sqrt{\frac{s_{1}^{2}}{n_{1}} + \frac{s_{2}^{2}}{n_{2}}}} > z_{\alpha} - \frac{\mu_{1} - \mu_{2} - \triangle_{l}}{\sqrt{\frac{s_{1}^{2}}{n_{1}} + \frac{s_{2}^{2}}{n_{2}}}}\right)
\end{align}
これより、
\begin{equation}
    -z_{\beta} = z_{\alpha} - \frac{\mu_{1} - \mu_{2} - \triangle_{l}}{\sqrt{\frac{s_{1}^{2}}{n_{1}} + \frac{s_{2}^{2}}{n_{2}}}}
\end{equation}
となるため、$n_{1}$について解くことでサンプルサイズの計算式を得る。
\begin{equation}
    n_{1} = \frac{(z_{\alpha/2} + z_{\beta})^{2} \left(s_{1}^{2} + \frac{s_{2}^{2}}{r}\right)}{(\mu_{1} - \mu_{2} - \triangle_{l})^{2}}
\end{equation}

帰無仮説が$\mu_{1} - \mu_{2} > \triangle_{u}$である場合も同様に導出できる。
\begin{align}
    1 - \beta & = P(T < -z_{\alpha})                                                                                                                                                                                                                                \\
              & = P\left(\frac{\bar{X}_{1} - \bar{X}_{2} - \triangle_{u}}{\sqrt{\frac{s_{1}^{2}}{n_{1}} + \frac{s_{2}^{2}}{n_{2}}}} < -z_{\alpha}\right)                                                                                                            \\
              & = P\left(\frac{\bar{X}_{1} - \bar{X}_{2} - (\mu_{1} - \mu_{2})}{\sqrt{\frac{s_{1}^{2}}{n_{1}} + \frac{s_{2}^{2}}{n_{2}}}} < -z_{\alpha} - \frac{\mu_{1} - \mu_{2} - \triangle_{u}}{\sqrt{\frac{s_{1}^{2}}{n_{1}} + \frac{s_{2}^{2}}{n_{2}}}}\right)
\end{align}
これより、
\begin{equation}
    z_{\beta} = -z_{\alpha} - \frac{\mu_{1} - \mu_{2} - \triangle_{u}}{\sqrt{\frac{s_{1}^{2}}{n_{1}} + \frac{s_{2}^{2}}{n_{2}}}}
\end{equation}
となるため、$n_{1}$について解くことでサンプルサイズの計算式を得る。
\begin{equation}
    n_{1} = \frac{(-z_{\alpha/2} - z_{\beta})^{2} \left(s_{1}^{2} + \frac{s_{2}^{2}}{r}\right)}{(\mu_{1} - \mu_{2} - \triangle_{u})^{2}} = \frac{(z_{\alpha/2} + z_{\beta})^{2} \left(s_{1}^{2} + \frac{s_{2}^{2}}{r}\right)}{(\mu_{1} - \mu_{2} - \triangle_{u})^{2}}
\end{equation}


\section{非劣性検定}

\subsection{どのような検定か}

非劣性検定は、ABテストにおいて現行バージョンを群1、新たなバージョンを群2とした時、群2が群1から劣化していないかを確かめるのに用いることができる。

2群の母平均$\mu_{1}, \mu_{2}$について、許容できる差を$\triangle_{l}$とするとき、帰無仮説$H_{0}$と対立仮説$H_{1}$は
\begin{gather}
    H_{0}: \mu_{1} - \mu_{2} < \triangle_{l} \\
    H_{1}: \triangle_{l} \le \mu_{1} - \mu_{2}
\end{gather}
とする。帰無仮説を棄却するとき、許容できる程度の劣化しかなかったと言える。

Welchのt検定を用いることができ、検定統計量は
\begin{equation}
    T_{l} = \frac{\bar{X}_{1} - \bar{X}_{2} - \triangle_{l}}{\sqrt{\frac{s_{1}^{2}}{n_{1}} + \frac{s_{2}^{2}}{n_{2}}}}
\end{equation}
である。棄却域は$T > t_{\alpha}(\nu)$である。

同等性検定のうち、lower equivalence boundに対する検定のみを行うと思えば良い。

\subsection{サンプルサイズの設計}

同等性検定における$\triangle_{l}$の時の計算式と同じである。


\section{マンホイットニーのU検定}

\subsubsection{どのような検定か}

「現代数理統計学」のセクション12.2に説明されている。値を元にランキングをつけて、順序を利用した検定統計量により検定を実施する。

\subsubsection{サンプルサイズの計算}

「Sample Size Determination for Some Common Nonparametric Tests」に書いてある。
ただし、この論文は片側検定の式なので、両側検定に使う場合は$z_{\alpha}$を$z_{\alpha/2}$に変える必要がある。


\section{多重比較の補正}

\subsection{Holm法}

\subsubsection{検定の方法}

多重比較における帰無仮説を$H_{0, i}\ (i = 1, \ldots, n)$とする。また、各検定におけるp値は$p_{i}$で、単調増加するように並んでいるものとする。$i = 1$から検定を開始する。$p_{i} < \frac{\alpha}{n + 1 - i}$ならば$H_{0, i}$を棄却し、$i + 1$に進める。$H_{0, i}$を棄却できない場合は検定を終了する。

\subsubsection{FWERを有意水準以下に抑えられることの証明}

標本空間$\Omega$を1回の実験で得られるデータ全体の集合（あるユーザーが何回CVしたなど）として、確率空間$(\Omega, \mathcal{F}, P)$を定義する。$I = \{i | H_{0, i}は真である\}$とし、$m = |I|$とする。

$i^{\ast} = \argmin_{i \in I} p_{i}$とする。$i^{\ast} \le n - m + 1$より$\frac{1}{n + 1 - i^{\ast}} \le \frac{1}{m}$を得る。FWERは、少なくとも1回帰無仮説が真である場合に誤って棄却する確率である。Holm法の検定方法より、少なくとも1回帰無仮説が真である場合に棄却されるためには、$H_{0, i^{\ast}}$が棄却される必要がある。$H_{0, i^{\ast}}$が棄却されなかった場合そこで検定が終了するため、他の帰無仮説も棄却されないからである。$H_{0, i^{\ast}}$が棄却される条件は$p_{i^{\ast}} < \frac{\alpha}{n + 1 - i^{\ast}} \le \frac{\alpha}{m}$である。これより、
\begin{equation}
    \text{FWER} \le P\left(\min_{i \in I} p_{i} \le \frac{\alpha}{m}\right)
\end{equation}
となる。さらに整理して、
\begin{equation}
    \text{FWER} \le P\left(\min_{i \in I} p_{i} \le \frac{\alpha}{m}\right) = P\left(\bigcup_{i \in I} \left\{p_{i} \le \frac{\alpha}{m}\right\}\right) \le \sum_{i \in I} P\left(\left\{p_{i} \le \frac{\alpha}{m}\right\}\right) \le m \frac{\alpha}{m} = \alpha
\end{equation}
となる。ここでは以下を用いた。
\begin{enumerate}
    \item 確率測度の劣加法性
    \item 帰無仮説が真である場合にp値は一様分布するため、$P\left(\left\{p_{i} \le \frac{\alpha}{m}\right\}\right) \le \frac{\alpha}{m}$であること
\end{enumerate}
帰無仮説が真である時にp値が一様分布することは、確率積分変換から導出することができる。「数理統計学の基礎」の命題2.19を参考にした。

\subsubsection{Bonferroni法に対する優位性}

Bonferroni法は、全ての帰無仮説に対して有意水準を$\alpha/n$とする。Holm法よりも各検定における有意水準が小さく、検出力に課題がある。Holm法はFWERを有意水準$\alpha$以下に抑えつつ、Bonferroni法よりも高い検出力で検定が可能である。


\section{その他}

\subsection{CUPED}

論文の証明の行間を埋める。
\begin{equation}
    \triangle_{cv} = \hat{Y}_{cv}^{(t)} - \hat{Y}_{cv}^{(c)}
\end{equation}
を整理すると、
\begin{align}
    \triangle_{cv} & = \hat{Y}_{cv}^{(t)} - \hat{Y}_{cv}^{(c)}                                                                                                   \\
                   & = (\bar{Y}^{(t)} - \theta \bar{X}^{(t)} + \theta \mathbb{E}[X^{(t)}]) - (\bar{Y}^{(c)} - \theta \bar{X}^{(c)} + \theta \mathbb{E}[X^{(c)}]) \\
                   & = (\bar{Y}^{(t)} - \theta \bar{X}^{(t)}) - (\bar{Y}^{(c)} - \theta \bar{X}^{(c)})                                                           \\
\end{align}
となる。以下より、$\triangle_{cv}$は$\triangle$に対する不偏推定量である。
\begin{align}
    \mathbb{E}[\triangle_{cv}] & = \mathbb{E}[(\bar{Y}^{(t)} - \theta \bar{X}^{(t)}) - (\bar{Y}^{(c)} - \theta \bar{X}^{(c)})] \\
                               & = \mathbb{E}[\bar{Y}^{(t)}] - \mathbb{E}[\bar{Y}^{(c)}]                                       \\
                               & = \mathbb{E}[\triangle]
\end{align}
$\triangle_{cv}$の分散を計算する。
\begin{align}
    \text{Var}[\triangle_{cv}] & = \text{Var}[(\bar{Y}^{(t)} - \theta \bar{X}^{(t)}) - (\bar{Y}^{(c)} - \theta \bar{X}^{(c)})]                       \\
                               & = \text{Var}[(\bar{Y}^{(t)} - \bar{Y}^{(c)}) - \theta (\bar{X}^{(t)} - \bar{X}^{(c)})]                              \\
                               & = \text{Var}[\triangle_{Y} - \theta \triangle_{X}]                                                                  \\
                               & = \text{Var}[\triangle_{Y}] + \theta^{2}\text{Var}[\triangle_{X}] - 2\theta\text{Cov}[\triangle_{Y}, \triangle_{X}]
\end{align}
第1項と第2項は次のように整理できる。
\begin{align}
    \text{Var}[\triangle_{Y}] & = \text{Var}[\bar{Y}^{(t)} - \bar{Y}^{(c)}]                                                                  \\
                              & = \text{Var}\left[\frac{1}{n_{t}}\sum_{i = 1}^{n_{t}}Y_{i} - \frac{1}{n_{c}}\sum_{i = 1}^{n_{c}}Y_{i}\right] \\
                              & = \frac{1}{n_{t}}\text{Var}[Y] + \frac{1}{n_{c}}\text{Var}[Y]                                                \\
                              & = \left(\frac{1}{n_{t}} + \frac{1}{n_{c}}\right)\text{Var}[Y]                                                \\
    \text{Var}[\triangle_{X}] & = \left(\frac{1}{n_{t}} + \frac{1}{n_{c}}\right)\text{Var}[X]
\end{align}
ここで、第3項を整理するために用いる式を導出する。
\begin{align}
    \text{Cov}[A - B, C - D]                                            & = \mathbb{E}[(A - B)(C - D)] - \mathbb{E}[A - B]\mathbb{E}[C - D]                                                                                                                           \\
                                                                        & = \mathbb{E}[AC - AD - BC + BD] - (\mathbb{E}[A] - \mathbb{E}[B])(\mathbb{E}[C] - \mathbb{E}[D])                                                                                            \\
                                                                        & = \mathbb{E}[AC] - \mathbb{E}[AD] - \mathbb{E}[BC] + \mathbb{E}[BD]                                                                                                                         \\
                                                                        & \quad - \mathbb{E}[A]\mathbb{E}[C] + \mathbb{E}[A]\mathbb{E}[D] + \mathbb{E}[B]\mathbb{E}[C] - \mathbb{E}[B]\mathbb{E}[D]                                                                   \\
                                                                        & = (\mathbb{E}[AC] - \mathbb{E}[A]\mathbb{E}[C]) - (\mathbb{E}[AD] - \mathbb{E}[A]\mathbb{E}[D])                                                                                             \\
                                                                        & \quad - (\mathbb{E}[BC] - \mathbb{E}[B]\mathbb{E}[C]) + (\mathbb{E}[BD] - \mathbb{E}[B]\mathbb{E}[D])                                                                                       \\
                                                                        & = \text{Cov}[A, C] - \text{Cov}[A, D] - \text{Cov}[B, C] + \text{Cov}[B, D]                                                                                                                 \\
    \text{Cov}[aX, bY]                                                  & = \mathbb{E}[aXbY] - \mathbb{E}[aX]\mathbb{bY}                                                                                                                                              \\
                                                                        & = ab(\mathbb{E}[XY] - \mathbb{E}[X]\mathbb{E}[Y])                                                                                                                                           \\
                                                                        & = ab\text{Cov}[X, Y]                                                                                                                                                                        \\
    \text{Cov}\left[\sum_{i = 1}^{n}X_{i}, \sum_{j = 1}^{m}Y_{j}\right] & = \mathbb{E}\left[\left(\sum_{i = 1}^{n}X_{i} - \mathbb{E}\left[\sum_{i = 1}^{n}X_{i}\right]\right)\left(\sum_{j = 1}^{m}Y_{j} - \mathbb{E}\left[\sum_{j = 1}^{m}Y_{j}\right]\right)\right] \\
                                                                        & = \mathbb{E}\left[\left(\sum_{i = 1}^{n}X_{i} - \sum_{i = 1}^{n}\mathbb{E}[X_{i}]\right)\left(\sum_{j = 1}^{m}Y_{j} - \sum_{j = 1}^{m}\mathbb{E}[Y_{j}]\right)\right]                       \\
                                                                        & = \mathbb{E}\left[\sum_{i = 1}^{n}\sum_{j = 1}^{m}(X_{i} - \mathbb{E}[X_{i}])(Y_{j} - \mathbb{E}[Y_{j}])\right]                                                                             \\
                                                                        & = \sum_{i = 1}^{n}\sum_{j = 1}^{m} \mathbb{E}\left[(X_{i} - \mathbb{E}[X_{i}])(Y_{j} - \mathbb{E}[Y_{j}])\right]                                                                            \\
                                                                        & = \sum_{i = 1}^{n}\sum_{j = 1}^{m} \text{Cov}[X_{i}, Y_{j}]
\end{align}
上記の式を用いて、第3項を整理する。
\begin{align}
    \text{Cov}[\triangle_{X}, \triangle_{Y}] & = \text{Cov}[\bar{Y}^{(t)} - \bar{Y}^{(c)}, \bar{X}^{(t)} - \bar{X}^{(c)}]                                                                                                  \\
                                             & = \text{Cov}[\bar{Y}^{(t)}, \bar{X}^{(t)}] - \text{Cov}[\bar{Y}^{(t)}, \bar{X}^{(c)}] - \text{Cov}[\bar{Y}^{(c)}, \bar{X}^{(t)}] + \text{Cov}[\bar{Y}^{(c)}, \bar{X}^{(c)}] \\
\end{align}
SUTVAより、第2項と第3項は0である。第1項と第4項は$(t), (c)$の添え字を省略して、次のように整理できる。
\begin{align}
    \text{Cov}[\bar{Y}, \bar{X}] & = \text{Cov}\left[\frac{1}{n}\sum_{i = 1}^{n}Y_{i}, \frac{1}{n}\sum_{j = 1}^{n}X_{j}\right] \\
                                 & = \frac{1}{n^{2}}\text{Cov}\left[\sum_{i = 1}^{n}Y_{i}, \sum_{j = 1}^{n}X_{j}\right]        \\
                                 & = \frac{1}{n^{2}} \sum_{i = 1}^{n}\sum_{j = 1}^{n} \text{Cov}[Y_{i}, X_{j}]                 \\
                                 & = \frac{1}{n^{2}} \sum_{i = 1}^{n} \text{Cov}[Y_{i}, X_{i}]                                 \\
                                 & = \frac{1}{n^{2}} n \text{Cov}[Y, X]                                                        \\
                                 & = \frac{1}{n} \text{Cov}[Y, X]                                                              \\
                                 & = \frac{1}{n} \text{Cov}[X, Y]
\end{align}
ここでもSUTVAより、$i \neq j \Rightarrow \text{Cov}[Y_{i}, X_{j}] = 0$であることを用いた。これより、
\begin{align}
    \text{Cov}[\triangle_{X}, \triangle_{Y}] & =  \frac{1}{n_{t}} \text{Cov}[X, Y] + \frac{1}{n_{c}} \text{Cov}[X, Y] \\
                                             & = \left(\frac{1}{n_{t}} + \frac{1}{n_{c}}\right)\text{Cov}[X, Y]
\end{align}
以上より、
\begin{align}
    \text{Var}[\triangle_{cv}] & = \text{Var}[\triangle_{Y}] + \theta^{2}\text{Var}[\triangle_{X}] - 2\theta\text{Cov}[\triangle_{Y}, \triangle_{X}]                                                                                             \\
                               & = \left(\frac{1}{n_{t}} + \frac{1}{n_{c}}\right)\text{Var}[Y] + \theta^{2} \left(\frac{1}{n_{t}} + \frac{1}{n_{c}}\right)\text{Var}[X] - 2\theta \left(\frac{1}{n_{t}} + \frac{1}{n_{c}}\right)\text{Cov}[X, Y] \\
                               & = \left(\frac{1}{n_{t}} + \frac{1}{n_{c}}\right)(\text{Var}[Y] + \theta^{2} \text{Var}[X] - 2\theta \text{Cov}[X, Y])                                                                                           \\
\end{align}
となる。
\begin{align}
    \frac{d}{d\theta} (\text{Var}[Y] + \theta^{2} \text{Var}[Y] - 2\theta \text{Cov}[X, Y]) & = 2\theta \text{Var}[X] - 2\text{Cov}[X, Y] = 0
\end{align}
より、$\theta = \text{Cov}[X, Y] / \text{Var}[X]$で$\text{Var}[\triangle_{cv}]$は最小となる。これを代入すると、
\begin{align}
    \text{Var}[\triangle_{cv}] & = \left(\frac{1}{n_{t}} + \frac{1}{n_{c}}\right)\left(\text{Var}[Y] + \frac{\text{Cov}[X, Y]^{2}}{\text{Var}[X]^{2}} \text{Var}[X] - 2 \frac{\text{Cov}[X, Y]}{\text{Var}[X]} \text{Cov}[X, Y]\right) \\
                               & = \left(\frac{1}{n_{t}} + \frac{1}{n_{c}}\right)\left(\text{Var}[Y] - \frac{\text{Cov}[X, Y]^{2}}{\text{Var}[X]}\right)                                                                               \\
                               & = \left(\frac{1}{n_{t}} + \frac{1}{n_{c}}\right)\text{Var}[Y]\left(1 - \frac{\text{Cov}[X, Y]^{2}}{\text{Var}[X]\text{Var}[Y]}\right)                                                                 \\
                               & = \left(\frac{1}{n_{t}} + \frac{1}{n_{c}}\right)\text{Var}[Y]\left(1 - \rho^{2}\right)
\end{align}
となる。一方、
\begin{align}
    \text{Var}[\triangle] & = \text{Var}[\bar{Y}^{(t)} - \bar{Y}^{(c)}]                   \\
                          & = \text{Var}[\bar{Y}^{(t)}] + \text{Var}[\bar{Y}^{(c)}]       \\
                          & = \left(\frac{1}{n_{t}} + \frac{1}{n_{c}}\right)\text{Var}[Y]
\end{align}
であるから、
\begin{equation}
    \text{Var}[\triangle_{cv}] = \text{Var}[\triangle](1 - \rho^{2})
\end{equation}
を得る。

% \subsection{有限母集団からの非復元抽出}
% 大きさ$N$の有限母集団を考え、各個体の特性値を$a_{1}, \ldots, a_{N}$とする。また、$X_{i}, i = 1, \ldots, n$を標本として抽出された観測値とする。非復元無作為抽出とは、任意の互いに異なる$i_{1}, \ldots, i_{n}$について
% \begin{equation}
%     P(X_{1} = a_{i_{1}}, \ldots, X_{n} = a_{i_{n}}) = \frac{1}{N(N - 1) \cdots (N - n + 1)}
% \end{equation}
% となるような標本抽出の方法である。

% また、無作為非復元抽出では上記のように$X_{1}, \ldots, X_{n}$を1個ずつ順に抜き出すと考えず、n個の個体を一度に抜き出すと考えることもでき、
% \begin{equation}
%     P(\{X_{1}, \ldots, X_{n}\} = \{a_{i_{1}}, \ldots, a_{i_{n}}\}) = 1 / \binom{N}{n}
% \end{equation}
% となる。

% ここで、このように一度に取り出した標本から、あらためて1個ずつ無作為に抜き出せば、母集団から1個ずつ順に抽出した時と同じ状況となる。特にi番目の観測値$X_{i}$の周辺分布は$X_{1}$の周辺分布と同じであり、$(X_{i}, X_{j})$の2次元の同時分布は$(X_{1}, X_{2})$の同時分布と同じである。これらを示す。

% まず、任意の$k \in \{i_{s}\}_{s = 1}^{n}$に対し、$X_{1} = a_{k}$となる確率は
% \begin{align}
%     P(X_{1} = a_{k})
%      & = \frac{{}_{N - 1} P_{n - 1}}{{}_{N} P_{n}}             \\
%      & = \frac{(N - 1)! / ((N - 1) - (n - 1))!}{N! / (N - n)!} \\
%      & = \frac{(N - 1)! / (N - n)!}{N! / (N - n)!}             \\
%      & = \frac{(N - 1)!}{N!}                                   \\
%      & = \frac{1}{N}
% \end{align}
% となるが、これは任意の$i = 1, \ldots, n$についても同様だから、$X_{i}$の周辺分布は$X_{1}$の周辺分布と同じである。

% 次に、任意の$k, l \in \{i_{s}\}_{s = 1}^{n}\ (k \neq l)$に対し、$X_{1} = a_{k}, X_{2} = a_{l}$となる確率は
% \begin{align}
%     P(X_{1} = a_{k}, X_{2} = a_{l})
%      & = \frac{{}_{N - 2} P_{n - 2}}{{}_{N} P_{n}}             \\
%      & = \frac{(N - 2)! / ((N - 2) - (n - 2))!}{N! / (N - n)!} \\
%      & = \frac{(N - 2)! / (N - n)!}{N! / (N - n)!}             \\
%      & = \frac{(N - 2)!}{N!}                                   \\
%      & = \frac{1}{N(N - 1)}
% \end{align}
% となるが、これは任意の$i, j \in \{1, \ldots, n\}\ (i \neq j)$についても同様だから、$X_{i}, X_{j}$の同時分布は$X_{1}, X_{2}$の同時分布と同じである。

% ここで、有限母集団の母平均$\mu$と母分散$\sigma^{2}$を次のように定義する。
% \begin{gather}
%     \mu = \frac{1}{N} \sum_{i = 1}^{N} a_{i} \\
%     \sigma^{2} = \frac{1}{N} \sum_{i = 1}^{N} (a_{i} - \mu)^{2}
% \end{gather}
% 標本平均$\bar{X} = \frac{1}{n} \sum_{i = 1}^{n} X_{i}$の期待値と分散を評価する。まず、期待値は
% \begin{align}
%     E[\bar{X}]
%      & = E\left[\frac{1}{n} \sum_{i = 1}^{n} X_{i}\right] \\
%      & = \frac{1}{n} \sum_{i = 1}^{n} E\left[X_{i}\right] \\
%      & = \frac{1}{n} \cdot n\mu                           \\
%      & = \mu
% \end{align}
% となる。次に分散は、
% \begin{align}
%     \text{Var}[\bar{X}]
%      & = \text{Var}\left[\frac{1}{n} \sum_{i = 1}^{n} X_{i}\right]                                                   \\
%      & = \sum_{i = 1}^{n} \frac{1}{n^{2}} \text{Var}[X_{i}] + 2\sum_{i < j} \frac{1}{n^{2}} \text{Cov}[X_{i}, X_{j}] \\
%      & = \frac{1}{n^{2}} \left(\sum_{i = 1}^{n} \text{Var}[X_{i}] + 2\sum_{i < j} \text{Cov}[X_{i}, X_{j}]\right)    \\
%      & = \frac{1}{n^{2}} \left(\sum_{i = 1}^{n} \text{Var}[X_{i}] + \sum_{i \neq j} \text{Cov}[X_{i}, X_{j}]\right)  \\
%      & = \frac{1}{n^{2}} \left(n\text{Var}[X_{1}] + n(n - 1)\text{Cov}[X_{1}, X_{2}]\right)
% \end{align}
% となる。$X_{i}$と$X_{1}$、$X_{i}, X_{j}$と$X_{1}, X_{2}$の周辺分布が同じことを利用し、分散および共分散をまとめた。

% ここで、
% \begin{align}
%     \sum_{i = 1}^{N} (a_{i} - \bar{a})^{2}
%      & = \sum_{i = 1}^{N} (a_{i}^{2} - 2a_{i}\bar{a} + \bar{a}^{2})                      \\
%      & = \sum_{i = 1}^{N} a_{i}^{2} - 2\bar{a} \sum_{i = 1}^{N}a_{i} + N\bar{a}^{2}      \\
%      & = \sum_{i = 1}^{N}a_{i}^{2} - 2\bar{a}N\bar{a} + N\bar{a}^{2}                     \\
%      & = \sum_{i = 1}^{N}a_{i}^{2} - 2N\bar{a}^{2} + N\bar{a}^{2}                        \\
%      & = \sum_{i = 1}^{N}a_{i}^{2} - N\bar{a}^{2}                                        \\
%      & = \sum_{i = 1}^{N}a_{i}^{2} - N\left(\frac{1}{N} \sum_{i = 1}^{N}a_{i}\right)^{2} \\
%      & = \sum_{i = 1}^{N}a_{i}^{2} - \frac{1}{N} \left(\sum_{i = 1}^{N}a_{i}\right)^{2}
% \end{align}
% 及び任意の$i \neq j$について$P(X_{1} = a_{i}, X_{2} = a_{j}) = 1 / N(N - 1)$であることに注意すると、$\text{Cov}[X_{1}, X_{2}]$は
% \begin{align}
%     \text{Cov}[X_{1}, X_{2}]
%      & = E[X_{1}X_{2}] - E[X_{1}]E[X_{2}]                                                                                                                               \\
%      & = \sum_{i \neq j}P(X_{1} = a_{i}, X_{2} = a_{j})a_{i}a_{j} - \left(\frac{1}{N}\sum_{i = 1}^{N}a_{i}\right)^{2}                                                   \\
%      & = \frac{1}{N(N - 1)} \sum_{i \neq j}a_{i}a_{j} - \left(\frac{1}{N}\sum_{i = 1}^{N}a_{i}\right)^{2}                                                               \\
%      & = \frac{1}{N(N - 1)}\left(\left(\sum_{i = 1}^{N} a_{i}\right)^{2} - \sum_{i = 1}^{N}a_{i}^{2}\right) - \left(\frac{1}{N}\sum_{i = 1}^{N}a_{i}\right)^{2}         \\
%      & = \frac{1}{N(N - 1)}\left(\sum_{i = 1}^{N}a_{i}\right)^{2} - \frac{1}{N(N - 1)}\sum_{i = 1}^{N}a_{i}^{2} - \frac{1}{N^{2}}\left(\sum_{i = 1}^{N}a_{i}\right)^{2} \\
%      & = \left(\frac{1}{N(N - 1)} - \frac{1}{N^{2}}\right)\left(\sum_{i = 1}^{N}a_{i}\right)^{2} - \frac{1}{N(N - 1)}\sum_{i = 1}^{N}a_{i}^{2}                          \\
%      & = \frac{N - (N - 1)}{N^{2}(N - 1)}\left(\sum_{i = 1}^{N}a_{i}\right)^{2} - \frac{1}{N(N - 1)}\sum_{i = 1}^{N}a_{i}^{2}                                           \\
%      & = \frac{1}{N^{2}(N - 1)}\left(\sum_{i = 1}^{N}a_{i}\right)^{2} - \frac{1}{N(N - 1)}\sum_{i = 1}^{N}a_{i}^{2}                                                     \\
%      & = -\frac{1}{N(N - 1)}\left(\sum_{i= 1}^{N}a_{i}^{2} - \frac{1}{N}\left(\sum_{i = 1}^{N}a_{i}\right)^{2}\right)                                                   \\
%      & = -\frac{1}{N(N - 1)}\sum_{i = 1}^{N}(a_{i} - \bar{a})^{2}                                                                                                       \\
%      & = -\frac{\sigma^{2}}{N - 1}
% \end{align}
% となる。従って$\text{Var}[\bar{X}]$は
% \begin{align}
%     \text{Var}[\bar{X}]
%      & = \frac{1}{n^{2}} \left(n\text{Var}[X_{1}] + n(n - 1)\text{Cov}[X_{1}, X_{2}]\right)        \\
%      & = \frac{1}{n^{2}} \left(n\sigma^{2} + n(n - 1)\left(-\frac{\sigma^{2}}{N - 1}\right)\right) \\
%      & = \frac{\sigma^{2}}{n} - \frac{n - 1}{n(N - 1)}\sigma^{2}                                   \\
%      & = \frac{\sigma^{2}}{n}\left(1 - \frac{n - 1}{N - 1}\right)                                  \\
%      & = \frac{\sigma^{2}}{n} \frac{(N - 1) - (n - 1)}{N - 1}                                      \\
%      & = \frac{\sigma^{2}}{n} \frac{N - n}{N - 1}
% \end{align}
% となる。まとめると、
% \begin{gather}
%     E[\bar{X}] = \mu \\
%     \text{Var}[\bar{X}] = \frac{N - n}{N - 1}\frac{\sigma^{2}}{n}
% \end{gather}
% である。


\renewcommand{\refname}{参考文献}

\begin{thebibliography}{99}
    \bibitem{tokyo} 東京大学教養学部統計学教室（編）『統計学入門 基礎統計学I』東京大学出版会、1991年。 \url{https://www.amazon.co.jp/dp/4130420658}
    \bibitem{kubokawa2017} 久保川達也『現代数理統計学の基礎』共立出版、2017年。\url{https://www.amazon.co.jp/dp/4320111664}
    \bibitem{takemura2020} 竹村彰通（2020）『新装改訂版 現代数理統計学』学術図書出版社.
    \bibitem{jss2023} 日本統計学会 編（2023）『増訂版 日本統計学会公式認定 統計検定1級対応 統計学』東京図書.
    \bibitem{shimizu2019} 清水泰隆『統計学への確率論、その先へ―ゼロからの測度論的理解と漸近理論への架け橋』内田老鶴圃、2019年。\url{https://www.amazon.co.jp/dp/B09HQGPM6L}
    \bibitem{nagata2003} 永田靖『サンプルサイズの決め方 (統計ライブラリー)』朝倉書店、2003年。\url{https://www.amazon.co.jp/dp/4254126654}
    \bibitem{kohavi2021} Ron Kohavi, Diane Tang, Ya Xu 著, 大杉直也 訳（2021）『A/Bテスト実践ガイド ――真のデータドリブンへ至る信頼できる実験とは』アスキー・ドワンゴ.
    \bibitem{lakens2017} Lakens, D. "Equivalence Tests: A Practical Primer for t Tests, Correlations, and Meta-Analyses." \textit{Social Psychological and Personality Science} 8.4 (2017): 355-362. \url{https://pmc.ncbi.nlm.nih.gov/articles/PMC5502906/pdf/10.1177_1948550617697177.pdf}
    \bibitem{ahn2012} Ahn, S., \& Park, S. H. "Understanding equivalence and noninferiority testing." \textit{Korean Journal of Pediatrics} 55.11 (2012): 403-405. \url{https://pmc.ncbi.nlm.nih.gov/articles/PMC3510268/pdf/kjped-55-403.pdf}
    \bibitem{noether1987} Gottfried E. Noether (1987) "Sample Size Determination for Some Common Nonparametric Tests", Journal of the American Statistical Association, 82(398), 645-647.
    \bibitem{deng2013} Deng, A., Xu, Y., Kohavi, R., \& Walker, T. "Improving the Sensitivity of Online Controlled Experiments by Utilizing Pre-Experiment Data." \textit{Proceedings of the Sixth ACM International Conference on Web Search and Data Mining} (2013): 123-132. \url{https://dl.acm.org/doi/pdf/10.1145/2433396.2433413}
    \bibitem{emuyn_equivalence} emuyn. "同等性検定." \url{https://www.emuyn.net/stats/lecture/%E5%90%8C%E7%AD%89%E6%80%A7%E6%A4%9C%E5%AE%9A}
    \bibitem{robinson_slutsky} Robinson, T. S. "10 Econometric Theorems: Slutsky's Theorem." \url{https://bookdown.org/ts_robinson1994/10EconometricTheorems/slutsky.html}
    \bibitem{korem2024} Allon Korem. "How to plan test duration when using CUPED." \textit{Statsig Blog} (2024). \url{https://www.statsig.com/blog/how-to-plan-test-duration-cuped}
\end{thebibliography}

\end{document}